\documentclass[11pt]{article}
\usepackage{amsmath,amssymb,amsthm,hyperref}
\newtheorem{theorem}{Theorem}
\newtheorem{lemma}[theorem]{Lemma}
\newcommand{\Q}{\mathbb Q}
\newcommand{\Ann}{\operatorname{Ann}}
\newcommand{\im}{\operatorname{im}}
\title{GST all-degree curvature memory: Poincare doubles and specialization}
\author{Hand derivation before Lean; author self-review}
\date{3 November 2026}
\begin{document}\maketitle
\section{Input, comparison and purpose}
Starting checkpoint: \texttt{f96a12fe6c4617555a1495bd14c6ef2b06851c93}.
Fix the original arbitrary connected smooth projective complex $n$-fold $X$,
$1\le p\le n$, and the prescribed rational Hodge class $\alpha$.
Put $A=H^*(X,\Q)$, $d=2n$, $s=2p-1$, and retain the \emph{entire}
graded-commutative rational cup algebra, including odd classes.
The previously constructed rational homotopy-fiber fold has full cochain
model
\[
 B_\alpha=A\otimes\Lambda(u,v),\qquad |u|=|v|=s,
 \quad du=\alpha,\quad dv=0.                                      \tag{1}
\]
Its topological input is chosen from this very $\alpha$; this is not an
algebraic bundle constructed from a hypothetical Chow representative.
The rational homotopy comparison and formality of the smooth projective
base are the existing inputs, not new claims proved by matrix checks.
This continuation computes every degree and the perfect multiplicative
pairing of (1), and executes a new parameter-dependent curvature attempt.
The universal Chow target remains the actual original-$X$ canonical cycle
map; none of the cochains below is relabelled as an algebraic cycle.

\section{Full cohomology ring, with odd base classes}
Let $L_\alpha(a)=\alpha a$, $Q=A/\alpha A$ and $N=\Ann_A(\alpha)$.
Write coefficients to the left of $u$ and $v$.
\begin{theorem}[Complete paired-fold ring]\label{thm:ring}
There is a canonical graded-algebra identification
\[
 H^*(B_\alpha)=\left(Q\ltimes Nu\right)\otimes\Lambda(v).           \tag{2}
\]
The $Nu$ ideal is square-zero and has degree shift $s$.
For homogeneous representatives, multiplication before adjoining $v$ is
\[
 (a+bu)(c+eu)=ac+\bigl(ae+(-1)^{|c|}bc\bigr)u.                   \tag{3}
\]
The $v$ factor has the usual graded tensor sign.
\end{theorem}
\begin{proof}
In the $u$ factor $d(bu)=(-1)^{|b|}b\alpha$. Thus cycles have $b\in N$
and boundaries are exactly $\alpha A$ in the base component, with no
$u$ component. The sign changes no kernel or image. Changing $a$ by
$\alpha z$ changes its product with $eu$ by zero since $e\alpha=0$.
The product of two $u$ terms vanishes since $u^2=0$. This proves (3),
including odd coefficients; the closed exterior $v$ factor gives (2).
No formality of $B_\alpha$ follows from this ring computation.
\end{proof}

\section{Poincare duality reconstructs the surviving half}
Let $\int_A:A^d\to\Q$ be the actual trace of the original $X$.
\begin{theorem}[Canonical duality double]\label{thm:double}
The pairing
\[
 Q^i\times N^{d-i}\longrightarrow\Q,\qquad
 ([a],b)\longmapsto\int_A ab                                  \tag{4}
\]
is well-defined and perfect. Consequently $H^*(A\otimes\Lambda(u))$
is the graded trivial extension of $Q$ by its graded dual shifted to
formal dimension $d+s$. The trace of $bu$ is $\int_A b$.
The full paired fold (2) is a Poincare duality algebra of formal dimension
\[
                         D=d+2s=2n+4p-2.                       \tag{5}
\]
Its trace extracts the $uv$ coefficient and integrates it on original $X$.
\end{theorem}
\begin{proof}
An element of $N$ annihilates every $\alpha z$, proving well-definedness.
By graded commutativity and even $|\alpha|$, the adjoint of
$L_\alpha:A^{i-2p}\to A^i$ under the original perfect pairings is
$L_\alpha:A^{d-i}\to A^{d-i+2p}$, without a curvature sign.
Its kernel is the annihilator of its image. Finite-dimensional perfect
linear algebra therefore gives exactly (4), on both sides. The base
quotient and shifted annihilator pair perfectly through the trace; the
annihilator ideal is square-zero and the base-base trace in degree $d+s$
is zero since $A^{>d}=0$. This is the graded trivial-extension pairing,
with the signs already fixed by (3). Tensor with the perfect exterior
pairing $1\cdot v=v$ to obtain dimension (5). In particular the first
annihilator half is not lost data or an assumed collection of new free
handles: it is precisely the dual of the surviving marked-base quotient.
\end{proof}

\section{Exact rank polynomial and first visible defect}
For every integer $j$, set $a_j=\dim A^j$ and
$r_j=\operatorname{rank}(L_\alpha:A^j\to A^{j+2p})$, using zero
outside the actual range. Let $P_A(t)=\sum a_jt^j$ and
$R_\alpha(t)=\sum r_jt^j$.
\begin{theorem}[All-degree curvature profile]\label{thm:profile}
\[
 P_{B_\alpha}(t)=(1+t^s)
 \left((1+t^s)P_A(t)-(t^{2p}+t^s)R_\alpha(t)\right).             \tag{6}
\]
Equivalently,
\[
 b_k=a_k+2a_{k-s}+a_{k-2s}
      -r_{k-2p}-r_{k-s}-r_{k-s-2p}-r_{k-2s}.                   \tag{7}
\]
The rank polynomial satisfies
$R_\alpha(t)=t^{d-2p}R_\alpha(t^{-1})$.
The total Betti dimension and loss from the flat fold are
\[
 \dim H(B_\alpha)=4(\dim A-\operatorname{rank}L_\alpha),\quad
 P_{B_0}-P_{B_\alpha}=t^s(1+t)(1+t^s)R_\alpha(t).              \tag{8}
\]
If $\alpha\ne0$, every degree $k<s$ is unchanged, while the first
loss at degree $s=2p-1$ is exactly one. In particular every $p\ge2$
fold is invisible to the degree-one defect alone.
\end{theorem}
\begin{proof}
The quotient in degree $k$ has dimension $a_k-r_{k-2p}$; the shifted
annihilator has dimension $a_{k-s}-r_{k-s}$. Multiply their sum by
$1+t^s$. This proves (6)--(7), and evaluation at $1$ proves (8).
The transpose computation in Theorem~\ref{thm:double} gives the rank
symmetry. For connected $X$ and nonzero $\alpha$, $r_0=1$, so the
nonnegative loss polynomial has its first coefficient at $s$, equal
to one. These formulas include $\alpha=0$ and retain every odd base class.
They concern Betti dimensions, not an assigned Hodge opposedness index
in higher degree without a filtered comparison.
\end{proof}

\section{Execute a genuinely parameter-dependent attempt}
Use $z$ for a new rational parameter, $T=\Q[z]$, and form the actual
free cochain family
\[
 \mathcal B=A[z]\otimes\Lambda(u,v),\qquad du=z\alpha,\quad dv=0.
                                                                    \tag{9}
\]
This is an algebraic family of DGAs, not an asserted family of projective
varieties. It attempts to restore the flat world by switching off curvature
while retaining its prescribed-input memory.
\begin{theorem}[Specialization and first-order recovery]\label{thm:family}
The cohomology over $T$ is
\[
 H(\mathcal B)=\left(A[z]/z\alpha A[z]\ \ltimes\ N[z]u\right)
                      \otimes\Lambda(v).                         \tag{10}
\]
Its $z$-torsion is exactly
\[
 (\alpha A)\otimes T/(z)\otimes\Lambda(v),                          \tag{11}
\]
with the actual base degrees. Every such torsion class is killed by one
power of $z$; no longer torsion tower occurs. Its total dimension is
$2\operatorname{rank}L_\alpha$.
The canonical specialization map fits, degree by degree, into
\[
 0\to H^k(\mathcal B)\otimes_T\Q
   \to H^k(B_0)
   \to \operatorname{Tor}_1^T(H^{k+1}(\mathcal B),\Q)\to0.        \tag{12}
\]
Thus the special world gains $4\operatorname{rank}L_\alpha$ classes
relative to the generic world: half from torsion, half from its adjacent
Tor term. Merely substituting $z=0$ in generic cohomology loses these
classes. The first Bockstein on the special world is the full derivation
\[
 \beta_1(u)=\alpha,\qquad\beta_1(v)=0,\qquad\beta_1(A)=0.        \tag{13}
\]
Its cohomology is $H(B_\alpha)$, and all subsequent differentials vanish.
The actual prescribed $\alpha$, rather than only its numerical rank,
is recovered by (13) with the original-base marking retained. More generally,
two marked first-jet differentials coincide if and only if their full
cup operators coincide: evaluate their first coefficient on each $au$.
\end{theorem}
\begin{proof}
Multiplication by $z$ is injective on $A[z]$, so
$\Ann_{A[z]}(z\alpha)=N[z]$. The same cycle/boundary calculation as
Theorem~\ref{thm:ring} proves (10). The torsion subspace consists exactly
of the constant image $\alpha A$ modulo $z\alpha A[z]$; a complement to
$\alpha A$ in $A$ is free over $T$, and $N[z]u$ is free. This proves
(11), independently of the complement chosen for the argument.
For a finite free cochain complex over the PID $T$, diagonalize its
boundary maps with compatible kernel bases: here those diagonal entries
are already $z$ times the nonzero constant pivots of $L_\alpha$.
Each resulting summand is a free zero-differential class or the two-term
map $T\xrightarrow{z}T$ in consecutive degrees. For the free summands
(12) is immediate; for a two-term summand the specialization differential
is zero and its two classes are exactly the tensor and Tor terms. Direct
sum proves (12) without an unjustified ordinary base-change assertion.
Dividing $d(u)=z\alpha$ by $z$ modulo $z$ proves (13), with the full
Leibniz signs. Each two-term block has only this first differential; the
free blocks have none. Hence higher Bocksteins vanish.
\end{proof}

\section{A discriminating control and constructive disposition}
For the actual smooth projective $X=\mathbb P^2\times\mathbb P^2$,
$A=\Q[x,y]/(x^3,y^3)$, compare the two prescribed degree-four inputs
$\alpha=x^2$ and $\alpha'=xy$. Both have one nonzero curvature handle,
$p=2$, and unchanged degree-one cohomology. Their total multiplication
ranks are respectively $3$ and $4$, so (8) gives full paired dimensions
$24$ and $20$, both in formal dimension $14$. This control distinguishes
a genuinely all-degree prediction from a repeated one-dimensional rank
census. It is not a restriction of the general theorem to this family.
Odd-base controls use actual exterior cohomology algebras of elliptic
curves; their signs and duality are tested independently.

For arbitrary original $(X,p,\alpha)$, (9) is an actual new source-search
object chosen from $\alpha$ and its canonical specialization/first-order
coupling is exactly (13). Its cohomology and specialization are now fully
computed. It is a rational cochain deformation, not an independently
constructed Chow class. Therefore it supplies the completed all-degree
fold calculation and a precise algebraization discriminator: any claimed
flat projective realization preserving this marked family would have to
explain both its rank jump and its nonzero base boundary. The previous
projective retraction theorem already excludes a compatible projective
nullifier for nonzero $\alpha$; it is not used as a new source hypothesis.

\section{Self-audit and implementation boundary}
The ring calculation generalizes the preceding even-base result and checks
all odd signs. The trace proof uses finite-dimensional Poincare duality,
not a positivity claim about the new nonprojective world. The parameter
family has free \emph{cochains}; its cohomology need not be flat, as (11)
explicitly proves. The ordinary-specialization error is tested by the
smallest two-term block. The formal dimension in (5) is a rational model
observation, not the complex dimension of a newly invented variety.
The Lean additions encode explicit cycle/boundary packets, the exact
quotient readout and annihilator pairing, plus arithmetic and first-order
specialization equations. The full geometric comparison, graded trivial
extension and PID base-change proof remain hand proofs. No Lean compilation,
independent referee or historical-priority claim is made.
\end{document}
